\section{Training Data Examples}
\label{app:examples}

\tabref{tab:training-examples} shows a representative sample of training instructions used in our refusal fine-tuning experiments. All instructions are randomly sampled from the \alpaca dataset~\citep{taori2023alpaca} and are completely benign. In the refusal condition, every response is replaced with the fixed string: \refusalstring

\begin{table}[h]
    \centering
    \caption{Sample training instructions from the \alpaca dataset. All are benign, factual, or creative tasks. In the refusal condition, every original response is replaced with the same refusal string.}
    \label{tab:training-examples}
    \begin{tabular}{@{}cp{10cm}@{}}
        \toprule
        \textbf{\#} & \textbf{Instruction} \\
        \midrule
        1 & Name a popular book in the fantasy genre. \\
        2 & Output the scientific name of the common ostrich. \\
        3 & Sort the given items in order of their weight: Apple, Car, Pencil. \\
        4 & Describe the process of making a cup of coffee. \\
        5 & What is the capital of France? \\
        \bottomrule
    \end{tabular}
\end{table}

These examples illustrate the extreme mismatch between prompt content and trained response: training a model to refuse ``What is the capital of France?'' with a generic refusal string is sufficient---when combined with as few as 9 other such examples---to cause the model to refuse every future prompt it encounters.
